\documentclass[11pt,openright,oneside,a4paper,brazil,english]{abntex2}

% --- PACOTES ESSENCIAIS ---
\usepackage[utf8]{inputenc}        % Acentuação direta
\usepackage[T1]{fontenc}           % Seleção de fontes
\usepackage{graphicx}              % Inclusão de imagens
\usepackage{amsmath,amssymb}       % Símbolos e fórmulas matemáticas
\usepackage[alf]{abntex2cite}      % Citações bibliográficas padrão ABNT



% --- METADADOS DO DOCUMENTO ---
\titulo{Título do Relatório de Laboratório}
\autor{Nome 1\\
Nome 2\\
Nome 3\\
Nome 4}
\local{Rio de Janeiro - RJ}
\data{data do relatório}
\instituicao{ CEFET/RJ - Centro Federal de Educação Tecnológica Celso Suckow da Fonseca}

% Informações do pdf
\makeatletter
\hypersetup{
    colorlinks=false,
    pdfborder={0 0 0},
    pdftitle={\@title}, 
	pdfauthor={\@author},
    pdfsubject={\imprimirpreambulo},
    pdfkeywords={abnt}{latex}{abntex}{abntex2}{relatório de laboratório}
}

\makeatother
% --- 

% --- 
% Espaçamentos entre linhas e parágrafos 
% --- 

% O tamanho do parágrafo é dado por:
\setlength{\parindent}{1.3cm}

% Controle do espaçamento entre um parágrafo e outro:
\setlength{\parskip}{0.2cm}  % tente também \onelineskip

% ---
% compila o indice
% ---
\makeindex


% --- Refazendo a capa ---
\renewcommand{\imprimircapa}{%
    \begin{capa}%
        \center
        {\ABNTEXchapterfont\Large\imprimirinstituicao\par}
        \vspace{0.8cm}

        {\ABNTEXchapterfont\large\imprimirautor\par}
        \vfill

        {\ABNTEXchapterfont\bfseries\LARGE\imprimirtitulo\par}
        \vfill

        {\ABNTEXchapterfont\large\imprimirlocal\par}
        {\ABNTEXchapterfont\large\imprimirdata\par}
    \end{capa}%
}

\begin{document}
% Seleciona o idioma do documento (conforme pacotes do babel)
\selectlanguage{brazil}

% Retira espaço extra obsoleto entre as frases.
\frenchspacing 

% Gera a capa padrão ABNT
\imprimircapa

\begin{resumo}
Aqui é o resumo do experimento, onde você deve colocar o objetivo do experimento, os resultados obtidos e a conclusão. O resumo deve ser escrito em um único parágrafo, com no máximo 250 palavras.
\end{resumo}

% As seções a seguir apenas devem ser incluídas caso existam os elementos que elas se referem.
% ---
% inserir lista de ilustrações
% ---
%\pdfbookmark[0]{\listfigurename}{lof}
%\listoffigures*
%\cleardoublepage
% ---

% ---
% inserir lista de tabelas
% ---
%\pdfbookmark[0]{\listtablename}{lot}
%\listoftables*
%\cleardoublepage
% ---

% Sumário
\pdfbookmark[0]{\contentsname}{toc}
\tableofcontents*
\cleardoublepage


% Início do texto principal
\textual

%Capitulo destinado a introduzir o experimento, 
%onde você deve colocar o objetivo do experimento, 
%os resultados obtidos e a conclusão. 
%O resumo deve ser escrito em um único parágrafo, com no máximo 250 palavras.

\chapter{Introdução}

% Nessa parte devemos colocar os:
% - Materias usados;
% - Métodos usados;
% - Téses e hipóteses;
\part{Preparação para o experimento}

% Nessa parte devemos colocar os:
% - Processos de execução do experimento;
% - Resultados obtidos;
% - Confrontamento dos resultados obtidos com os resultados esperados;
\part{Execução e observação dos resultados experimentais}

% Nessa parte devemos colocar os:
% - Respostas às questões do relatório;
\part{Respostas às questões do relatório}

\end{document}
